\begin{definition}[Algebraic torus] \label{definition:algebraic_torus_over_a_field}
    An \hldef{(algebraic) torus over a field $k$} is an \CrefAndHyperrefIfExist{definition:algebraic_group_scheme_over_a_scheme}{algebraic group} $T$ over $k$ such that the base change $T_{\bark} = T \times_{k} \bark$ is isomorphic as an algebraic group over $\bark$ to \CrefAndHyperrefIfExist{definition:multiplicative_group_scheme_over_a_scheme}{$(\mathbb{G}_m/\bark)^n$} for some nonnegative integer $n$. In particular, any $l$ of $k$ such that $T \times_{k} l$ is isomorphic to $(\bbG_{m}/l)^n$ may be called a \hldef{splitting field of the torus}. There always exists a minimal splitting field of the  torus and it is a finite extension of $k$.
    Equivalently, an algebraic torus over $k$ is an \CrefAndHyperrefIfExist{definition:algebraic_torus_over_a_scheme}{algebraic torus} over \CrefAndHyperrefIfExist{definition:affine_scheme}{$\Spec k$}.

\end{definition}
